% !TeX root = ../main.tex

\begin{itemize}
	\item Je vais reprendre la notion d’UX design ici et à nouveau la travailler sur ce sujet spécifiquement. L’idée étant de se baser sur les questions qui ont été soulevées lors de l’établissement de la spécification de cette fonctionnalité. 
	% \item Méfiance utilisateurs et implémentation de l'utilisation : comment le gérer ? 
	% \item  Parcours utilisateurs anciens et nouveaux, quelle position adopter ? Forcer le changement ? Proposer la nouvelle solution comme une alternative ? L’insérer dans les parcours classiques identifier ? 
	
	\item \textbf{Plan :}
	\subitem 1. Je ne redéfinis pas l'UX Design, mais je pars de la conclusion de ma partie 2, chapitre 1, section 3 sur l'UX design appliqué au NLP, et remobilise les mêmes connaissances et sources. J'annonce dès le départ cet axe avec l'orientation vers l'indexation semi-automatisée dans l'application. 
	\subitem 2. Contrairement au NLP où on propose d'approfondir des manières déjà existantes d'utiliser l'application ou une nouvelle manière d'usage mais sans interférence sur la gestion des collections \textit{("chatbot")}, ici il s'agit de changer complètement la manière de travailler. Pas juste d'usage. Se pose la question de comment utiliser l'UX design pour faciliter le changement dans la manière de travailler en plus de faciliter l'usage de l'outil ? 
	\subitem 3. Reprise des enjeux vus sur le NLP mais appliqués à cet outil. Comment ne pas trop impacter l'utilisateur négativement tout en créant une réelle valeur ajoutée ? Question du changement de parcours utilisateurs : changement total ou partiel ?. Montrer toutes les difficultés qu'une telle fonctionnalité amène. L'idée est de jongler avec des concepts d'UX design pour montrer les difficultés d'implémentations qu'il risque d'y avoir et ce que ça peut impacter. 
	
\end{itemize}

Nous ne reviendrons pas dans cette partie sur les notions d'\gls{ux} et d'\gls{ui} largement traitée page \pageref{uxui}. Il s'agit ici de reprendre les considérations précédentes et de se questionner sur l'inclusion de l'indexation semi-automatisée dans un \gls{sgc}\newline

Nous avons donc mis en exergue la complexité que représente l'intégration du \gls{NLP}, sous deux outils différents, dans l'application. Cela revient à changer presque totalement le rapport de l'utilisateur avec son système de production par l'introduction de technologies envers lesquelles une certaines méfiances existent. Ce changement important ne doit pas provoquer de bouleversement graphiques et d'expérience utilisateur : si à l'échelle technique les changements sont important, l'impact utilisateur se doit d'être contrôlé. 

Finalement, le c\oe{ur} de l'enjeu réside dans l'adoption de ses outils, ou du moins leurs non rejet, ce qui passe par deux besoins. D'abords la transparence des processus et résultats. Ensuite, le besoin de placer l'humain au c\oe{ur} des processus notamment par la possibilité de valider ou invalider des résultats. \newline

Mais qu'en est-il de l'indexation semi-automatisée ? Les enjeux d'\gls{ux} et d'\gls{ui} sont-ils similaires ? Ont-ils les mêmes réponses ? La réponse est en réalité assez nuancée. 

Si \gls{NLP} change le rapport de l'utilisateur au système, cette proposition amène à changer entièrement une manière de travailler avec le logiciel SkinSoft et, à plus large échelle, au sein des collections d'\glspl{imgani}. On retrouve bien évidemment cette question de l'adoption de l'outil, qui trouve forcément les mêmes réponses.